% GENERATED FILE. Edit canonical lessons, then run npm run generate:books.
% canonical-source-hash: fnv1a64:4fddf434a6fe172f
% canonical-lessons: GU-C368-nasib, GU-C368-bhagya, GU-C368-sanyog, GU-C368-phalvu, GU-C368-nivadvu

\chapter{Luck and Fate}
\label{ch:gu-luck-and-fate}
\hlchaptermodality{368}

\begin{chapteropening}
\textbf{By the end of this chapter:} \emph{I can say what luck has to do with it.}

Complete the last lesson of chapter 368: I can say what luck has to do with it.
\end{chapteropening}

\section[\texorpdfstring{\gu{નસીબ} (nasīb)}{નસીબ (nasīb)}]{\gu{નસીબ} (nasīb) — luck}

\label{lesson:GU-C368-nasib}



\begin{quote}
Before the new one: say the Gujarati for an assumption, then the Gujarati for to assess.
\end{quote}

\subsection*{You'll want to know: \gu{નસીબ}}
\textbf{\gu{નસીબ}} — \emph{nasīb} — \textquotedblleft{}luck\textquotedblright{}.

\begin{quote}
\gu{જો} \gu{નસીબ} \gu{સારું} \gu{હશે} \gu{તો} \gu{આપણે} \gu{જીતીશું}.
\end{quote}

\begin{quote}
\emph{jo nasīb sārũ haśe to āpaṇe jītīśũ.}
\end{quote}

\begin{quote}
If our luck holds, we will win.
\end{quote}

A neuter noun.

\begin{cousinweb}
\textbf{\gu{નસીબ}} (\emph{nasīb}) is usually traced to Arabic \emph{naṣīb}, \textquotedblleft{}a share, a portion\textquotedblright{}, through Persian: the share life deals out.
\end{cousinweb}

\begin{grammarlens}[title={What you've built}]
One more word for saying what happens if.
\end{grammarlens}

\subsection*{Your turn}
\begin{itemize}
\raggedright
  \item \emph{Say it:} \emph{nasīb}
  \item \emph{Say it:} \emph{nasīb}, once more
  \item \emph{Say it:} \emph{jo nasīb sārũ haśe to āpaṇe jītīśũ.} — If our luck holds, we will win.
  \item \emph{Recall:} say the Gujarati for an assumption, then the Gujarati for to assess, then say \emph{nasīb} again
\end{itemize}

\begin{tcolorbox}[breakable,colback=teal!4,colframe=teal!35!black,title={Before you move on}]
What does \gu{નસીબ} mean? (\textquotedblleft{}Luck\textquotedblright{}.) Say it once more.
\end{tcolorbox}
\section[\texorpdfstring{\gu{ભાગ્ય} (bhāgya)}{ભાગ્ય (bhāgya)}]{\gu{ભાગ્ય} (bhāgya) — destiny}

\label{lesson:GU-C368-bhagya}



\begin{quote}
Before the new one: say the Gujarati for to assess, then the Gujarati for luck.
\end{quote}

\subsection*{You'll want to know: \gu{ભાગ્ય}}
\textbf{\gu{ભાગ્ય}} — \emph{bhāgya} — \textquotedblleft{}destiny\textquotedblright{}.

\begin{quote}
\gu{જો} \gu{ભાગ્ય} \gu{સારું} \gu{હોત} \gu{તો} \gu{તે} \gu{આજે} \gu{અમીર} \gu{હોત}.
\end{quote}

\begin{quote}
\emph{jo bhāgya sārũ hot to te āje amīr hot.}
\end{quote}

\begin{quote}
If fate had been kind, he would be rich today.
\end{quote}

A neuter noun.

\begin{cousinweb}
\textbf{\gu{ભાગ્ય}} (\emph{bhāgya}) is a tatsama, borrowed from Sanskrit \emph{bhāgya}, \textquotedblleft{}fortune, fate\textquotedblright{} (Monier-Williams).
\end{cousinweb}

\begin{grammarlens}[title={What you've built}]
One more word for saying what happens if.
\end{grammarlens}

\subsection*{Your turn}
\begin{itemize}
\raggedright
  \item \emph{Say it:} \emph{bhāgya}
  \item \emph{Say it:} \emph{bhāgya}, once more
  \item \emph{Say it:} \emph{jo bhāgya sārũ hot to te āje amīr hot.} — If fate had been kind, he would be rich today.
  \item \emph{Recall:} say the Gujarati for to assess, then the Gujarati for luck, then say \emph{bhāgya} again
\end{itemize}

\begin{tcolorbox}[breakable,colback=teal!4,colframe=teal!35!black,title={Before you move on}]
What does \gu{ભાગ્ય} mean? (\textquotedblleft{}Destiny\textquotedblright{}.) Say it once more.
\end{tcolorbox}
\section[\texorpdfstring{\gu{સંયોગ} (sanyog)}{સંયોગ (sanyog)}]{\gu{સંયોગ} (sanyog) — a coincidence}

\label{lesson:GU-C368-sanyog}



\begin{quote}
Before the new one: say the Gujarati for luck, then the Gujarati for destiny.
\end{quote}

\subsection*{You'll want to know: \gu{સંયોગ}}
\textbf{\gu{સંયોગ}} — \emph{sanyog} — \textquotedblleft{}a coincidence\textquotedblright{}.

\begin{quote}
\gu{જો} \gu{આ} \gu{સંયોગ} \gu{ન} \gu{હોત} \gu{તો} \gu{અમે} \gu{એક} \gu{જ} \gu{ટ્રેનમાં} \gu{ન} \gu{હોત}.
\end{quote}

\begin{quote}
\emph{jo ā sanyog na hot to ame ek ja ṭrenmā̃ na hot.}
\end{quote}

\begin{quote}
But for this coincidence, we would not have been on the same train.
\end{quote}

A masculine noun.

\begin{cousinweb}
\textbf{\gu{સંયોગ}} (\emph{sanyog}) is a tatsama, borrowed from Sanskrit \emph{saṃyoga}, \textquotedblleft{}a joining, a connection\textquotedblright{} (Monier-Williams). \textbf{\gu{સંજોગ}} (\emph{sanjog}), \textquotedblleft{}a circumstance\textquotedblright{}, is the same word in an older, inherited shape.
\end{cousinweb}

\begin{grammarlens}[title={What you've built}]
One more word for saying what happens if.
\end{grammarlens}

\subsection*{Your turn}
\begin{itemize}
\raggedright
  \item \emph{Say it:} \emph{sanyog}
  \item \emph{Say it:} \emph{sanyog}, once more
  \item \emph{Say it:} \emph{jo ā sanyog na hot to ame ek ja ṭrenmā̃ na hot.} — But for this coincidence, we would not have been on the same train.
  \item \emph{Recall:} say the Gujarati for luck, then the Gujarati for destiny, then say \emph{sanyog} again
\end{itemize}

\begin{tcolorbox}[breakable,colback=teal!4,colframe=teal!35!black,title={Before you move on}]
What does \gu{સંયોગ} mean? (\textquotedblleft{}A coincidence\textquotedblright{}.) Say it once more.
\end{tcolorbox}
\section[\texorpdfstring{\gu{ફળવું} (phaḷvũ)}{ફળવું (phaḷvũ)}]{\gu{ફળવું} (phaḷvũ) — to bear fruit}

\label{lesson:GU-C368-phalvu}



\begin{quote}
Before the new one: say the Gujarati for destiny, then the Gujarati for a coincidence.
\end{quote}

\subsection*{You'll want to know: \gu{ફળવું}}
\textbf{\gu{ફળવું}} — \emph{phaḷvũ} — \textquotedblleft{}to bear fruit\textquotedblright{}.

\begin{quote}
\gu{જો} \gu{તમે} \gu{સતત} \gu{મહેનત} \gu{કરશો} \gu{તો} \gu{તે} \gu{ફળશે}.
\end{quote}

\begin{quote}
\emph{jo tame satat mahenat karśo to te phaḷśe.}
\end{quote}

\begin{quote}
If you keep working hard, it will pay off.
\end{quote}

Past (it …): \gu{તે} \gu{કેરીનો} \gu{છોડ} \gu{ફળ્યો}, \emph{te kerīno choḍ phaḷyo}.

\begin{cousinweb}
\textbf{\gu{ફળવું}} (\emph{phaḷvũ}) is inherited from Sanskrit \emph{phálati}, \textquotedblleft{}bears fruit\textquotedblright{}; CDIAL lists the Gujarati verb (CDIAL 9057).
\end{cousinweb}

\begin{grammarlens}[title={What you've built}]
One more word for saying what happens if.
\end{grammarlens}

\subsection*{Your turn}
\begin{itemize}
\raggedright
  \item \emph{Say it:} \emph{phaḷvũ}
  \item \emph{Say it:} \emph{phaḷvũ}, once more
  \item \emph{Say it:} \emph{jo tame satat mahenat karśo to te phaḷśe.} — If you keep working hard, it will pay off.
  \item \emph{Recall:} say the Gujarati for destiny, then the Gujarati for a coincidence, then say \emph{phaḷvũ} again
\end{itemize}

\begin{tcolorbox}[breakable,colback=teal!4,colframe=teal!35!black,title={Before you move on}]
What does \gu{ફળવું} mean? (\textquotedblleft{}To bear fruit\textquotedblright{}.) Say it once more.
\end{tcolorbox}
\section[\texorpdfstring{\gu{નીવડવું} (nīvaḍvũ)}{નીવડવું (nīvaḍvũ)}]{\gu{નીવડવું} (nīvaḍvũ) — to turn out}

\label{lesson:GU-C368-nivadvu}



\begin{quote}
Before the new one: say the Gujarati for a coincidence, then the Gujarati for to bear fruit.
\end{quote}

\subsection*{You'll want to know: \gu{નીવડવું}}
\textbf{\gu{નીવડવું}} — \emph{nīvaḍvũ} — \textquotedblleft{}to turn out\textquotedblright{}.

\begin{quote}
\gu{જો} \gu{આ} \gu{દવા} \gu{સારી} \gu{નીવડશે} \gu{તો} \gu{ઘણા} \gu{દર્દીઓ} \gu{સાજા} \gu{થશે}.
\end{quote}

\begin{quote}
\emph{jo ā davā sārī nīvaḍśe to ghaṇā dardīo sājā thaśe.}
\end{quote}

\begin{quote}
If this medicine turns out well, many patients will get better.
\end{quote}

Past (it …): \gu{તે} \gu{વાત} \gu{સાચી} \gu{નીવડી}, \emph{te vāt sācī nīvḍī}.

\begin{cousinweb}
\textbf{\gu{નીવડવું}} (\emph{nīvaḍvũ}) CDIAL lists the Gujarati verb, \textquotedblleft{}to turn out, to come up\textquotedblright{}, under a reconstructed \emph{nirvarati}, \textquotedblleft{}restrains\textquotedblright{} (CDIAL 7392); how \textquotedblleft{}restrain\textquotedblright{} became \textquotedblleft{}turn out\textquotedblright{} is not explained there.
\end{cousinweb}

\begin{grammarlens}[title={What you've built}]
One more word for saying what happens if.
\end{grammarlens}

\subsection*{Your turn}
\begin{itemize}
\raggedright
  \item \emph{Say it:} \emph{nīvaḍvũ}
  \item \emph{Say it:} \emph{nīvaḍvũ}, once more
  \item \emph{Say it:} \emph{jo ā davā sārī nīvaḍśe to ghaṇā dardīo sājā thaśe.} — If this medicine turns out well, many patients will get better.
  \item \emph{Recall:} say the Gujarati for a coincidence, then the Gujarati for to bear fruit, then say \emph{nīvaḍvũ} again
\end{itemize}

\begin{tcolorbox}[breakable,colback=teal!4,colframe=teal!35!black,title={Before you move on}]
What does \gu{નીવડવું} mean? (\textquotedblleft{}To turn out\textquotedblright{}.) Say it once more.
\end{tcolorbox}
